\section{Model Assumptions and Nomenclature}
\subsection{Fundamental Assumptions}

To develop a mathematically tractable model that effectively captures the key dynamics of the battery, this study introduces the following fundamental assumptions in establishing the second-order RC equivalent circuit model, based on the physical characteristics of lithium-ion batteries and typical application scenarios:

\begin{table}[htbp]
\centering
\caption{Fundamental Model Assumptions}
\label{tab:assumptions}
\begin{tabular}{p{0.22\textwidth}p{0.73\textwidth}}
\toprule
\textbf{Category} & \textbf{Description} \\
\midrule
\textbf{1. Circuit Equivalence} & The complex internal electrochemical processes (e.g., lithium-ion intercalation/deintercalation, solid-phase diffusion, charge transfer) can be sufficiently approximated by a linear circuit network composed of lumped-parameter elements (resistors, capacitors, voltage sources). Specifically, a topology consisting of an open-circuit voltage source in series with an ohmic resistor and two parallel RC branches is adopted. \\
\textbf{2. Electrochemical Simplifications} & a) The internal temperature of the battery is assumed uniform and constant, neglecting dynamic effects of temperature variation on model parameters.\\
& b) The open-circuit voltage $V_{\text{oc}}$ is solely a single-valued function of SOC, and its hysteresis effect is considered negligible.\\
& c) The polarization dynamics can be described by first-order linear kinetics, with time constants remaining constant during a single discharge process.\\
\textbf{3. Operational \& Boundary Conditions} & a) The primary focus is on constant-power ($P_0$) discharge conditions, where the output power is precisely set by the external load and remains constant.\\
& b) Discharge starts from a fully relaxed battery state, i.e., with zero initial polarization voltages: $V_{p1}(0) = V_{p2}(0) = 0$.\\
& c) The Coulombic efficiency $\eta$ is constant during discharge, typically assumed to be 1.\\
& d) The power conversion efficiency from the battery terminals to the load, $\eta_{\text{conv}}$, is constant.\\
\textbf{4. Parameter Properties} & a) For a single simulation, the polarization resistances $R_1, R_2$ and capacitances $C_1, C_2$ are treated as constants, independent of SOC, current, or temperature. Their values are taken from the open dataset from the University of Maryland\cite{Calce2026}\\
& b) The nominal battery capacity $C_n$ and the State of Health (SOH) are known constants.\\
& c) The battery terminal voltage must not fall below a preset cutoff voltage $V_{\text{min}}$; discharge is terminated when this limit is reached.\\
\bottomrule
\end{tabular}
\end{table}

The aforementioned assumptions aim to strike a balance between model complexity and practicality. 
Assumptions 1 and 2 form the physical basis of the second-order RC equivalent circuit, simplifying the distributed-parameter, 
nonlinear electrochemical system into a lumped-parameter, linear dynamic system. Assumption 3 defines the core operational 
condition (constant-power discharge) and initial/boundary conditions for this study. Assumption 4 significantly simplifies 
the parameter identification and model-solving process, which is common in engineering practice. Under this combined framework, 
the dynamic equations are established and can be solved efficiently using numerical methods. 
All subsequent simulations and analyses are conducted within this framework.

\begin{figure}[H] 
    \centering
    \includegraphics[width=0.8\textwidth]{flowchart_battery_model.png}
    \caption{Framework flowchart of smartphone battery discharge modeling}
    \label{fig:flowchart}
\end{figure}

\subsection{Nomenclature}
The main symbols used in this paper and their descriptions are listed in Table \ref{tab:symbols}.
\begin{table}[H]
\centering
\caption{Nomenclature}
\label{tab:symbols}
\begin{tabular}{p{0.18\textwidth}p{0.65\textwidth}p{0.1\textwidth}}
\toprule
\textbf{Symbol} & \textbf{Description} & \textbf{Unit} \\
\midrule
$V_{\text{oc}}$ & Battery open-circuit voltage, function of SOC & V \\
SOC & State of Charge, $0 \leq \text{SOC} \leq 1$ & -- \\
SOH & State of Health, indicating capacity fade, $0 < \text{SOH} \leq 1$ & -- \\
$V_t$ & Battery terminal voltage & V \\
$I$ & Battery current (positive for discharge) & A \\
$P_0$ & Constant discharge power & W \\
$R_0$ & Ohmic internal resistance & $\Omega$ \\
$R_1, C_1$ & Polarization resistance and capacitance of the 1st RC network (describing electrochemical polarization) & $\Omega$, F \\
$R_2, C_2$ & Polarization resistance and capacitance of the 2nd RC network (describing concentration polarization) & $\Omega$, F \\
$V_{p1}$ & Polarization voltage across the 1st RC network & V \\
$V_{p2}$ & Polarization voltage across the 2nd RC network & V \\
$\tau_1 = R_1C_1$ & Time constant of the 1st RC network & s \\
$\tau_2 = R_2C_2$ & Time constant of the 2nd RC network & s \\
$C_n$ & Nominal battery capacity & Ah \\
$\eta$ & Coulombic efficiency (typically 1 for discharge) & -- \\
$\eta_{\text{conv}}$ & Power conversion efficiency & -- \\
$V_{\text{min}}$ & Discharge cutoff voltage & V \\
$\text{SOC}_{\text{min}}$ & Minimum allowable SOC & -- \\
$t$ & Time & s \\
$\Delta t$ & Numerical integration time step & s \\
$\mathbf{x} = [\text{SOC}, V_{p1}, V_{p2}]^\top$ & System state vector & -- \\
$P_{\text{max}}$ & Instantaneous maximum deliverable power & W \\
$Q_{\text{disch}}$ & Cumulative discharged capacity & Ah \\
$I_{\text{avg}}$ & Average discharge current & A \\
$t_{\text{end}}$ & Discharge termination time & s \\
\bottomrule
\end{tabular}
\end{table}